\documentclass[12pt]{extarticle}

% ============================================================
% PACKAGES
% ============================================================

% Encodage et langue
\usepackage[french]{babel}

% Mise en page
\usepackage[a4paper, margin=2cm]{geometry}

\usepackage[francais,bloc,completemulti]{automultiplechoice}    

% Mathématiques
\usepackage{amsmath}
\usepackage{amssymb}

% Dessins
\usepackage{tikz}

% Listes
\usepackage{enumitem}

% Boîtes
\usepackage[most]{tcolorbox}

% Images
\usepackage{graphicx}

% Mise en page avancée
\usepackage{fancyhdr}
\usepackage{multicol}

% Tableaux
\usepackage{booktabs}
\usepackage{tabularx}
\usepackage{colortbl}

\usepackage[normalem]{ulem}

\usepackage{pgfplots}
\pgfplotsset{compat=1.18}
\usepackage{pgf-pie}


\usepackage{ifthen}

% ============================================================
% COULEURS
% ============================================================

% ------------------------------------------------------------
% Couleurs principales
% ------------------------------------------------------------

\definecolor{Bleu}{RGB}{40, 100, 180}
\definecolor{Vert}{RGB}{50, 150, 90}
\definecolor{Orange}{RGB}{230, 140, 40}

\definecolor{BleuClair}{RGB}{235, 245, 255}
\definecolor{VertClair}{RGB}{235, 250, 240}
\definecolor{OrangeClair}{RGB}{255, 245, 225}

% ------------------------------------------------------------
% Couleurs complémentaires
% ------------------------------------------------------------

% Rappels de cours — violet
\definecolor{Violet}{RGB}{125, 85, 170}
\definecolor{VioletClair}{RGB}{245, 238, 252}

% Théorèmes — rouge
\definecolor{Rouge}{RGB}{190, 55, 55}
\definecolor{RougeClair}{RGB}{255, 238, 238}

% Lemme — indigo
\definecolor{Indigo}{RGB}{75, 85, 155}
\definecolor{IndigoClair}{RGB}{240, 241, 252}

% Points méthode — turquoise
\definecolor{Turquoise}{RGB}{35, 145, 150}
\definecolor{TurquoiseClair}{RGB}{235, 250, 250}

% Erreurs / attention — jaune
\definecolor{Jaune}{RGB}{210, 165, 35}
\definecolor{JauneClair}{RGB}{255, 250, 225}

% ------------------------------------------------------------
% Neutres
% ------------------------------------------------------------

\definecolor{GrisClair}{RGB}{245, 245, 245}
\definecolor{Gris}{RGB}{100, 100, 100}


% ============================================================
% COMMANDES PERSONNALISÉES
% ============================================================

\newcommand{\bb}[1]{\textbf{#1}}
\newcommand{\ii}[1]{\textit{#1}}

\newcommand{\cm}{\,\mathrm{cm}}
\newcommand{\m}{\,\mathrm{m}}

\newcommand{\reponse}[1]{\boxed{#1}}

% ============================================================
% TABLEAUX
% ============================================================

\newcommand{\ligneentete}{%
    \rowcolor{Bleu}%
    \color{white}\bfseries
}

\newcommand{\entete}[1]{%
    \cellcolor{Bleu}\color{white}\bb{#1}%
}


\NewDocumentCommand{\hl}{O{JauneClair}m}{%
  \tcbox[
    on line,
    colback=#1,
    colframe=#1,
    boxrule=0pt,
    arc=4pt,
    left=0pt,
    right=0pt,
    top=1pt,
    bottom=1pt
  ]{#2}%
}

\NewDocumentCommand{\hlb}{O{JauneClair}m}{%
  \tcbox[
    on line,
    colback=#1,
    colframe=#1,
    boxrule=0pt,
    arc=4pt,
    left=0pt,
    right=0pt,
    top=1pt,
    bottom=1pt
    ]{\bb{#2}}%
}
\newlength{\answerwidth}


\newcommand{\answer}[2][VertClair]{%
    \ifmmode
        \settowidth{\answerwidth}{$\displaystyle #2$}%
        \ifanswers
            {\setlength{\fboxsep}{0pt}%
             \colorbox{#1}{$\displaystyle #2$}}%
        \else
            \makebox[0pt][l]{%
                \makebox[\answerwidth][l]{%
                    \leaders\hbox to 0.5em{\hss.\hss}\hfill
                }%
            }%
            \phantom{#2}%
        \fi
    \else
        \settowidth{\answerwidth}{#2}%
        \ifanswers
            \hl[#1]{#2}%
        \else
            \makebox[\answerwidth][l]{%
                \leaders\hbox to 0.5em{\hss.\hss}\hfill
            }%
        \fi
    \fi
}


\newcommand{\graphanswer}[1]{%
    \ifanswers
        #1%
    \else
        0%
    \fi
}

\newcommand{\exo}[1]{%
    \includegraphics[
        width=\linewidth,
        keepaspectratio
    ]{#1}%
}

\newcommand{\nodata}{}
\newcommand{\correctanswer}[1]{
    \explain{
            \newline\vspace{10pt}
            #1
    }
}

\newlength{\mylen}
\newtcolorbox{exercicebox}[1][\unskip]{
    colback=white,
    colframe=Bleu,
    boxrule=1pt,
    arc=3mm,
    title=\bb{Exercice #1},
    left=2mm,
    right=2mm,
    top=2mm,
    bottom=1mm
}

\newcounter{theorem}
\newtcolorbox{theoremebox}[1][\unskip]{
    colback=RougeClair,
    colframe=Rouge,
    boxrule=1pt,
    arc=3mm,
    title=\bb{Théorème \stepcounter{theorem}\thetheorem}\hspace{5pt}#1,
    left=5mm,
    right=5mm,
    top=4mm,
    bottom=4mm
}

\newcounter{probleme}
\newtcolorbox{probbox}[1][\unskip]{
    colback=RougeClair,
    colframe=Rouge,
    boxrule=1pt,
    arc=3mm,
    title=\bb{Problème #1},
    left=5mm,
    right=5mm,
    top=2mm,
    bottom=2mm
}

\newenvironment{tableau}[1]
{%
    \renewcommand{\arraystretch}{1.3}%
    \begin{center}%
    \begin{tabularx}{\linewidth}{#1}%
    \toprule
}
{%
    \bottomrule
    \end{tabularx}%
    \end{center}%
    \renewcommand{\arraystretch}{1}%
}



% Set page layout after the config files
\geometry{
  paper=a4paper,
  hmargin=1cm,
  headheight=1cm,
  headsep=.3cm,
  footskip=0cm,
  top=1.5cm,
  left=1.3cm,
  right=1.3cm,
  bottom=1.5cm
}

%\geometry{hmargin=1cm,headheight=1cm,headsep=.3cm,footskip=0cm,top=1.5cm,bottom=1.5cm}

\def\AMCotextReserved{\emph{(Reservé)}}
\AMCopenOpts{boxframerule=0pt,boxframerulecol=red,framerule=0pt}
\newdimen\AMCpostOquest\AMCpostOquest=-7mm
\AMCformHSpace=.3em
\AMCinterBquest=1ex

\makeatletter
\def\jepasse{%
\def\AMCbeginQuestion##1##2{\ifAMC@catalog\textbf{##1}\fi}
\AMCquestionNumberfalse}
\makeatother


\begin{document}

\exemplaire{30}{    

% ---- Header ------
\fbox{
    \begin{minipage}{.40\linewidth}
        \namefield{Nom: }
        Prénom :\\
        Date:
    \end{minipage}
}
\hfill
\begin{minipage}{.48\linewidth}
    \AMCboxStyle{size=1.7ex,down=.2ex}
    \AMCcodeGridInt[vertical=false]{etu}{2}
\end{minipage}
% ---- Fin Header ------


\newcommand{\hpt}[1]{
    \AMCboxStyle{shape=square,color=black,size=1.8ex,down=.2ex}
    \begin{minipage}{2ex}
        \jepasse
        \begin{question}{##1}\bareme{b=1}
            \begin{reponseshoriz}
                \bonne{}
            \end{reponseshoriz}
        \end{question}
    \end{minipage} 
}
\newcommand{\pt}[1]{
    \AMCboxStyle{shape=oval,color=black,size=2ex,down=.2ex}
    \begin{minipage}{2ex}
        \jepasse
        \begin{question}{##1}\bareme{b=1}
            \begin{reponseshoriz}
                \bonne{}
            \end{reponseshoriz}
        \end{question}
    \end{minipage} 
}
 
 \begin{exercicebox}[\it{(Priorités opératoires)}
     \hfill\pt{op1}\pt{op2}\pt{op3}\pt{op4}\pt{op5}\hpt{op5.5}\hpt{op6}]
    Effectuer les calculs suivants \bb{en détaillant les étapes}\\
    \begin{minipage}{0.48\textwidth}
        \begin{align*}
        &\text{A = }23-7+2 \\
        &\text{B = }2+5\times4 \\
        &\text{C = }6\times4\div2 
        \end{align*}
    \end{minipage}
    \hfill
    \begin{minipage}{0.48\textwidth}
        \begin{align*}
            &\text{D = }10-4+10\div2\\
            &\text{E = }7\times[10-(8-3)]\\
            &\text{F = } 2\times3\div4\times(18-16)
        \end{align*}
    \end{minipage}
\end{exercicebox}


\begin{exercicebox}[\it{(Vocabulaire des opérations)}
     \hfill\pt{voc1}\pt{voc2}\hpt{voc2.5}\hpt{voc3}
    ]
    \begin{enumerate}
        \item Traduis chacune des expressions suivantes sous forme de phrases:
        \begin{align*}
            &\text{A = }10+6 \hspace{2cm}\text{B = }12-4\hspace{2cm}\text{C =
            }21\times4\hspace{2cm}\text{D = }64\div(3+5)
        \end{align*}
    \item Écris en une seule expression \bb{puis calcule} :
        \begin{enumerate}
            \item le produit de 8 par 6
            \item la différence entre 13 et le produit de 2 par 3
        \end{enumerate}
    \end{enumerate}
\end{exercicebox}


\begin{exercicebox}[\it{(Puissances et fractions)}
     \hfill\pt{pf1}\pt{pf2}\pt{pf3}\pt{pf4}\pt{pf5}\hpt{pf5.5}\hpt{pf6}]
    Effectuer les calculs suivants \bb{en détaillant les étapes}\\
    \begin{minipage}{0.48\textwidth}
        \begin{align*}
            &\text{A = }\dfrac{2\times18-6}{2+3} \\
            &\text{B = } 3+2^2
        \end{align*}
    \end{minipage}
    \hfill
    \begin{minipage}{0.48\textwidth}
        \begin{align*}
            &\text{C = }(14-2\times4)^2-3^3 \\
            &\text{D = }\dfrac{10\times(12-7)^2}{2\times(23-18)^2} 
        \end{align*}
    \end{minipage}
\end{exercicebox}


\begin{probbox}[
     \hfill\pt{prob1}\pt{prob2}\pt{prob3}\pt{prob4}\hpt{prob4.5}\hpt{prob5}]
    Un collège organise une sortie : le car coûte 480 €, l'entrée du musée 7,50 € par élève et 12 € par accompagnateur. Il y a 54 élèves et 6 accompagnateurs. Le foyer verse une aide de 150 €.
    \begin{enumerate}
        \item Écris, en une seule expression, le coût total de la sortie après
            l'aide du foyer. 
        \item Calcule ce coût total.
        \item Quelle somme chaque élève doit-il payer si le coût restant est 
            partagé équitablement entre les élèves ? 
        \item Le collège annonce "15 € par élève ". Cette annonce est-elle correcte ?
    \end{enumerate}
\end{probbox}

%\begin{probbox}
%    Un groupe d'amis décident de faire une excursion dont le prix est fixé par
%    une agence de voyage.
%
%    \begin{itemize}
%        \item 1er cas : si chacun des amis verse 150€, il manque alors 1500€ pour
%            l'excursion.
%        \item 2eme cas : si chacun des amis vers 200€, il y a 200€ de trop.
%    \end{itemize}
%
%    Quel est le nombre de participant à l'excursion ?\\
%    Quel est le prix fixé par l'agence de voyage pour l'excursion ?
%\end{probbox}

}   
\end{document}
